% Negative Trials and Null Results — Section for ch21-clinical-trials.tex
% Generated from reviewer feedback action plan A3
% To integrate: \input this file at the end of ch21-clinical-trials.tex

\section{Negative Trials and Null Results}
\label{sec:negative-trials}

Understanding why treatments fail is as informative as understanding why they succeed. This section documents major negative, null, and inconclusive clinical trials in ME/CFS. Each entry identifies not only the outcome but the lessons that should guide future trial design.

\begin{limitation}[Publication Bias in Negative Trials]
Negative trials are systematically underreported in the medical literature. The trials documented here represent those that were published; an unknown number of negative results may remain in file drawers. This section necessarily reflects survivorship bias in published negative results.
\end{limitation}

\subsection{Immunological Interventions}
\label{subsec:neg-immunological}

\subsubsection{Rituximab: From Promise to Definitive Negative}

The rituximab story is the most instructive negative trial in ME/CFS research, illustrating how small-sample enthusiasm can collapse in larger, properly controlled studies.

\paragraph{Phase II (Fluge et al., 2011).}
A small double-blind, placebo-controlled trial ($n = 30$) of two rituximab infusions (500~mg/m$^2$) showed a response rate of 67\% in the rituximab group vs.\ 13\% placebo at 12 months, with responses typically appearing 3--8 months after infusion~\cite{Fluge2011rituximab}. The delayed response pattern suggested an autoimmune mechanism requiring time for pathogenic antibody turnover.

\paragraph{Phase III --- RituxME (Fluge et al., 2019).}
The definitive phase~III trial enrolled 151 patients across four Norwegian hospitals in a double-blind, placebo-controlled design with six rituximab infusions over 12 months~\cite{Fluge2019}.

\textbf{Primary outcomes}: Neither co-primary endpoint---self-reported fatigue (Fatigue Severity Scale) nor overall function (SF-36 Physical Function)---showed significant differences between groups ($p = 0.81$ for overall response rate: 26\% rituximab vs.\ 35\% placebo). The trial was unambiguously negative.

\paragraph{Lessons learned.}
\begin{itemize}
    \item \textbf{Heterogeneity dilutes subgroup signals}: A 67\% response in a small, possibly enriched sample collapsed to no signal in an unstratified larger trial.
    \item \textbf{B-cell depletion may target the wrong compartment}: Subsequent daratumumab (anti-CD38, targeting plasma cells) showed a 60\% response rate~\cite{Fluge2025daratumumab}, suggesting long-lived plasma cells---not B cells---are the relevant autoantibody source.
    \item \textbf{Biomarker stratification is essential}: Future immunotherapy trials require pre-selection by autoantibody status.
\end{itemize}

\subsubsection{Intravenous Immunoglobulin (IVIG): Contradictory Results}

IVIG trials produced frustratingly inconsistent results spanning clearly positive to clearly negative outcomes.

\paragraph{Lloyd et al.\ (1990).}
The first placebo-controlled IVIG trial ($n = 49$) enrolled patients with immunological abnormalities. Approximately 43\% of IVIG-treated patients showed significant functional improvement~\cite{Lloyd1990}. Low baseline CD4 count predicted response.

\paragraph{Vollmer-Conna et al.\ (1997).}
A larger follow-up ($n = 99$) tested three IVIG doses monthly for three months~\cite{VollmerConna1997}. No dose showed benefit over placebo on any outcome.

\paragraph{Lessons learned.}
\begin{itemize}
    \item \textbf{Patient selection determines outcome}: Lloyd's positive trial required immunological abnormalities; Vollmer-Conna's did not.
    \item \textbf{Immune biomarkers may predict response}: Low CD4 count, IgG subclass deficiency, and age at onset may identify IVIG-responsive subgroups.
\end{itemize}

\subsection{Behavioral and Rehabilitation Interventions}
\label{subsec:neg-behavioral}

\subsubsection{The PACE Trial: GET and CBT}

The PACE trial remains the most controversial study in ME/CFS research.

\paragraph{Original trial (White et al., 2011).}
A large ($n = 641$), four-arm, assessor-blind RCT tested CBT, graded exercise therapy (GET), adaptive pacing therapy, and specialist medical care alone~\cite{White2011pace}. The original publication reported CBT and GET were ``moderately effective,'' with recovery rates of approximately 22\%.

\paragraph{Reanalysis (Wilshire et al., 2018).}
Using the originally published protocol thresholds, Wilshire et al.\ found~\cite{Wilshire2018}:
\begin{itemize}
    \item Recovery rates dropped from 22\% to approximately 7--8\%
    \item No objective measures (6-minute walk, actigraphy, employment) improved
    \item ``The modest treatment effects do not exceed what could be accounted for by participant reporting biases''
\end{itemize}

In 2021, NICE removed GET and CBT-as-treatment from its ME/CFS guidelines.

\paragraph{Lessons learned.}
\begin{itemize}
    \item \textbf{Objective outcomes are non-negotiable}: Self-reported outcomes in unblinded behavioral trials cannot be sole primary endpoints.
    \item \textbf{Protocol fidelity matters}: Mid-trial changes to outcome thresholds undermine validity.
    \item \textbf{Data transparency is essential}: The reanalysis was possible only through Freedom of Information litigation.
\end{itemize}

\subsection{Antiviral Interventions}
\label{subsec:neg-antiviral}

\subsubsection{Acyclovir (Straus et al., 1988)}

An early RCT ($n = 27$) tested IV then oral acyclovir for persistent EBV~\cite{Straus1988}. No significant difference between groups.

\textbf{Lesson}: Wrong antiviral spectrum (acyclovir has minimal HHV-6 activity) and no virus-specific patient selection.

\subsubsection{Valganciclovir --- Mixed Results (Montoya et al., 2013)}

The EVOLVE trial ($n = 30$) of valganciclovir in patients with elevated HHV-6/EBV titers~\cite{Montoya2013valganciclovir}. Primary composite endpoint: not significant. However, individual endpoints showed significant improvements in mental fatigue ($p = 0.039$), fatigue severity ($p = 0.006$), and cognition ($p = 0.025$).

\textbf{Lesson}: Composite endpoints can mask real subset effects. Virus-selected populations show stronger signals.

\subsubsection{Rintatolimod (Ampligen): Regulatory Failure}

Phase~III ($n = 234$) showed marginal overall improvement in exercise tolerance~\cite{Strayer2020}. Post-hoc analysis found 51.2\% improvement in patients with 2--8 year disease duration~\cite{Strayer2020rintatolimod}. FDA advisory committee voted 8--5 against approval, citing reliance on post-hoc subgroup analyses.

\textbf{Lesson}: Post-hoc subgroup analyses cannot rescue a trial; the disease-duration effect requires prospective confirmation.

\subsection{Neurological and Cholinergic Interventions}
\label{subsec:neg-neurological}

\subsubsection{Galantamine (Blacker et al., 2004)}

The largest ME/CFS pharmacological trial ($n = 434$) tested galantamine at four doses across 35 centers~\cite{Blacker2004}. No statistically significant differences on any outcome measure.

\textbf{Lesson}: Definitive null result rules out peripheral cholinergic enhancement. Cognitive impairment in ME/CFS is not cholinergic.

\subsection{Antidepressant Interventions}
\label{subsec:neg-antidepressant}

\subsubsection{Fluoxetine (Vercoulen et al., 1996)}

Double-blind RCT ($n = 96$) enrolling both depressed and non-depressed CFS patients~\cite{Vercoulen1996}. Fluoxetine showed no effect on any CFS dimension in either subgroup.

\textbf{Lesson}: CFS is not depression. Serotonergic modulation does not address core fatigue.

\subsection{Endocrine Interventions}
\label{subsec:neg-endocrine}

\subsubsection{Hydrocortisone}

Two RCTs tested hydrocortisone replacement based on mild hypocortisolism findings. McKenzie et al.\ ($n = 70$): modest improvement offset by adrenal suppression~\cite{McKenzie1998}. Cleare et al.\ ($n = 32$): 28\% vs.\ 9\% normalized fatigue scores, but adrenal suppression even at low doses~\cite{Cleare1999}.

\textbf{Lesson}: Mild hypocortisolism may be a \textit{consequence} of ME/CFS rather than a cause, explaining why correction does not resolve symptoms.

\subsection{Summary Table}
\label{subsec:neg-summary-table}

\begin{table}[htbp]
\centering
\caption{Major negative, null, and inconclusive trials in ME/CFS.}
\label{tab:negative-trials}
\small
\begin{tabular}{@{}p{2.5cm}p{1.2cm}p{0.8cm}p{1.3cm}p{5.5cm}@{}}
\toprule
\textbf{Intervention} & \textbf{Design} & $n$ & \textbf{Result} & \textbf{Key Lesson} \\
\midrule
Rituximab III~\cite{Fluge2019} & DB-RCT & 151 & Negative & No benefit unstratified; subgroup may exist \\
\addlinespace
PACE (GET/CBT)~\cite{White2011pace} & AB-RCT & 641 & Contested & No objective benefit; reanalysis $\sim$7\%~\cite{Wilshire2018} \\
\addlinespace
Rintatolimod~\cite{Strayer2020} & DB-RCT & 234 & Mixed & Marginal overall; subset signal at 2--8~yr \\
\addlinespace
Galantamine~\cite{Blacker2004} & DB-RCT & 434 & Negative & Definitive null; not cholinergic \\
\addlinespace
Acyclovir~\cite{Straus1988} & DB-RCT & 27 & Negative & Wrong antiviral spectrum \\
\addlinespace
Valganciclovir~\cite{Montoya2013valganciclovir} & DB-RCT & 30 & Mixed & Primary negative; individual endpoints positive \\
\addlinespace
IVIG~\cite{VollmerConna1997} & DB-RCT & 99 & Negative & No benefit without immune pre-selection \\
\addlinespace
Fluoxetine~\cite{Vercoulen1996} & DB-RCT & 96 & Negative & CFS $\neq$ depression \\
\addlinespace
Hydrocortisone~\cite{McKenzie1998} & DB-RCT & 70 & Negative & Adrenal suppression outweighs benefit \\
\bottomrule
\multicolumn{5}{@{}l}{\footnotesize DB = double-blind; AB = assessor-blind; RCT = randomized controlled trial.}
\end{tabular}
\end{table}

\subsection{Cross-Cutting Lessons}
\label{subsec:neg-lessons}

\begin{enumerate}
    \item \textbf{ME/CFS is not one disease for treatment purposes.} Every trial enrolling unstratified patients produced null or weak results. Every trial pre-selecting by biomarker found stronger signals.

    \item \textbf{Subjective outcomes in unblinded trials are unreliable.} Future trials must include objective measures: actigraphy, two-day CPET, employment status, or validated functional capacity assessments.

    \item \textbf{Composite primary endpoints can mask real effects.} Selecting a single, clinically meaningful primary endpoint is preferable.

    \item \textbf{Post-hoc subgroup analyses are hypothesis-generating, not confirmatory.} Adaptive trial designs can incorporate pre-specified biomarker stratification.

    \item \textbf{The null hypothesis is informative.} Each negative trial narrows the hypothesis space: galantamine rules out cholinergic mechanisms; fluoxetine argues against serotonergic deficit; hydrocortisone reframes hypocortisolism as consequence not cause.
\end{enumerate}
